% EGA V, English translation by P. Blass, 1987 (cote 24): archive cote 24, site page 15 (PDF page 16).
% typescript. Draft compared: none.
% AI-transcribed by Claude Opus 5.5 (claude-opus-5-5, adaptive thinking, effort high), Message Batches API, batch ega24b (msgbatch_01RGJJeHdR96U16ZWFuUjtDx), 2026-11-13; not yet checked by eye.
\scanpage{24/15}
\mspage{15}

Let us note that in 2.10 and 2.11 we do not talk at all about
the corresponding geometric properties; the results anticipated
are of 'absolute' nature.  We now examine the properties of
geometric nature. (One could if one wanted to to take the
opportunity to change the n$^\circ$)

\uline{Theorem 2.12}

Suppose that $f{:}X \to \underline{P}$ is unramified.  Let $y \in Y_{\eta}$, $x$ its image
in X.  In order for X to be smooth over k at x it is ncessary and
sufficient that $Y_{\eta}$ should be smooth over $k(\eta)$ at y.

We may assume that k is algebraically closed.  If $Y_{\eta}$ is smooth
over $k(\eta)$ at y it is regular, thus since $Y_{\eta}$ is flat over X, X is
regular at x (reference), therefore it is smooth over k at x since
k is algebraically closed and thus perfect (reference).

For the converse we can (after replacing X by an open
neighborhood of x) assume that X is smooth, and (due to the
jacobian criterion of smoothness) to be defined in an open subset
U of $\underline{P}$ by $\underline{p}$ equations as $X = V(f_1,\ldots,f_{\underline{p}})$, where the differentials
$df_1,\ldots,df_{\underline{p}}$ are everywhere linearly independent.  By introducing
the affine coordinates $S_1,\ldots,S_n$ in $\underline{P}$v and the affine coordinates
$T_1$, $T_2,\ldots,T_n$ in a neighborhood of x (by choosing a hyperplane $H\infty$
[at infinity] not containing x) $y \to U_{k(\eta)}$ is then equal to
$V(f_1,\ldots,f_{\underline{p}}, \Sigma\, S_iT_i - 1)$ and it suffices to verify that the
differentials (relative to $k(\eta)$) of $f_1,\ldots,f_{\underline{p}}$, $(\Sigma\, S_iT_i) - 1)$ are
linearly independent.  However these differentials are nothing
else but the sections of $(\Omega 1_{u/k}) \otimes_k k(\quad)$ [illegible]* as follows:
$df_1,\ldots,df_{\underline{p}}$, $\Sigma\, S_i dT_i$.  Since the $df_i$ are linearly independent of

\margin{* Ask A.G.}
